\documentclass[12pt, a4paper]{article}
\usepackage{amsmath, amssymb, geometry}
\usepackage{graphicx}
\usepackage{xcolor}
\geometry{margin=1in}

\title{\textbf{Dynamic Rule of Thirds\\Heatmap Mathematics}}
\author{Multi-Actor FPV Visualization Algorithm}
\date{}

\begin{document}

\maketitle

\section*{1. Mathematical Objective}
The goal of the algorithm is to construct a differentiable, dynamic "cost-reward" landscape across the drone's 2D First Person View (FPV) camera plane. The landscape is designed to perfectly map to the photographic "Rule of Thirds", assigning the highest mathematical reward (heat) when an actor is positioned precisely on one of the four grid intersection points (the power points). It is additionally weighted dynamically to favor the primary actor over the secondary actor.

\section*{2. FPV Plane and Power Point Coordinates}
The FPV camera uses a standard perspective projection with focal length $f$ and a predefined sensor bounds vector, mapping the horizontal physical scale as $vs$ and the vertical physical scale using an aspect ratio $\rho$. 

The camera's viewing plane is bounded horizontally by $\left[-vs, +vs\right]$ and vertically by $\left[-vs \cdot \rho, +vs \cdot \rho\right]$. By definition, the Rule of Thirds grid intersects exactly at horizontal and vertical one-third marks. Thus, our four power points, $\mathbf{P}_i = (cu_i, cv_i)$ for $i \in \{1,2,3,4\}$, are natively defined as:
\begin{align*}
\mathbf{P}_1 &= \left(-\frac{vs}{3}, -\frac{vs\cdot\rho}{3}\right) & \mathbf{P}_2 &= \left(+\frac{vs}{3}, -\frac{vs\cdot\rho}{3}\right) \\
\mathbf{P}_3 &= \left(-\frac{vs}{3}, +\frac{vs\cdot\rho}{3}\right) & \mathbf{P}_4 &= \left(+\frac{vs}{3}, +\frac{vs\cdot\rho}{3}\right)
\end{align*}

\section*{3. Distance-Based Dynamic Activation}
To make the heatmap react naturally to the actors entering the scene, we calculate the continuous Euclidean distance between the actors' projected FPV coordinates and each power point. Let the projected coordinate of the primary actor be $\mathbf{A}^{(1)} = (u_1, v_1)$ and the secondary actor be $\mathbf{A}^{(2)} = (u_2, v_2)$.

Each power point $i$ is assigned a dynamic activation scalar $A_i$. The activation rests at a baseline minimum ($\beta = 0.15$) so the grid is visible globally. As an actor approaches the power point, the activation increases exponentially. The primary actor carries a full priority multiplier $W_1 = 1.0$, while the secondary actor carries a reduced threshold $W_2 = 0.4$.

The activation score $A_i$ for the $i$-th power point is calculated as:
\begin{equation}
A_i = \min\left(1.15, \;\beta + W_1 e^{-\frac{\|\mathbf{A}^{(1)} - \mathbf{P}_i\|^2}{2\sigma_A^2}} + W_2 e^{-\frac{\|\mathbf{A}^{(2)} - \mathbf{P}_i\|^2}{2\sigma_A^2}}\right)
\end{equation}
Where $\sigma_A = 0.15$ defines the physical sensor drop-off radius for maximum activation.

\section*{4. The Final Gaussian Sum Heatmap}
With the four power point positions $\mathbf{P}_i$ and their respective real-time activations $A_i$ established, we compute the final heat signature $H(u,v)$ for any arbitrary microscopic pixel coordinate $(u,v)$ on the camera screen. 

We apply a 2D Gaussian Mixture Model (sum of 2D Gaussians) spread uniformly using standard deviations $\sigma_u = \frac{vs}{4.5}$ and $\sigma_v = \frac{vs\cdot\rho}{4.5}$. 

The final intensity mathematical equation mapping the color of each pixel $(u, v)$ is:
\begin{equation}
H(u, v) = \sum_{i=1}^{4} A_i \cdot \exp\left( - \frac{(u - cu_i)^2}{2\sigma_u^2} - \frac{(v - cv_i)^2}{2\sigma_v^2} \right)
\end{equation}

When rendered by the simulation natively, $H$ is mapped continuously to an \texttt{:inferno} or \texttt{:cgrad} color gradient overlay. The mathematical result effectively allows the Model Predictive Control solver to track the continuous differential sum toward these peak values to achieve photographic framing.

\end{document}
